\section{Introduction}

Cognition arises from coordinated activity across many interacting brain regions, yet most of what we understand about cortical computation comes from studies of single areas. The orientation columns of primary visual cortex and the entorhinal grid-cell map are iconic examples of local-circuit function, but neither speaks directly to how populations distributed across the cortex cooperate to support behaviour. Advances in large-scale electrophysiology and population imaging have begun to close this gap. High-density silicon probes now enable simultaneous recording from hundreds of sites in behaving rodents~\cite{jun2017neuropixels}, and synchronised multi-probe experiments reveal that choices, actions, and engagement are coded across far-flung cortical and subcortical regions~\cite{steinmetz2019distributed}. In parallel, cortex-wide calcium imaging has shown that single-trial neural dynamics are shaped by richly varied movements distributed across many areas~\cite{musall2019single,stringer2019spontaneous}. Together, these tools have opened a new era of research into the distributed basis of cognition.

Working memory is a paradigmatic distributed function. It is the brain's ability to hold information online across short gaps between sensory input and behaviour~\cite{baddeley2012working}, and delay-period persistent activity has been reported in many prefrontal and parietal regions of primates for decades~\cite{funahashi1989mnemonic,goldmanrakic1995cellular}. More recent work, summarised in reviews of the distributed nature of working memory~\cite{christophel2017distributed,leavitt2017sustained}, argues that mnemonic codes are not the property of a single area but emerge from interactions among many. In mice, inactivation of anterolateral motor cortex and of thalamic partners disrupts short-term memory and upcoming-action signals~\cite{guo2017maintenance,inagaki2019discrete,gilad2018behavioral}, while mediodorsal thalamus sustains prefrontal activity during maintenance~\cite{bolkan2017thalamic,schmitt2017thalamic}. These findings place distributed, recurrent activity at the centre of mammalian working memory.

A concurrent line of work has built connectome-based computational models to explain this distributed code. In the macaque, multiregional models endowed with a macroscopic gradient of synaptic excitation and dopamine modulation reproduce hierarchical timescales, graded persistent activity, and prefrontal dominance~\cite{murray2014hierarchy,mejias2022mechanisms,froudistwalsh2021dopamine}. The shared premise is that, in primates, spine counts increase along the cortical hierarchy and the resulting recurrent excitation gradient is what allows association cortex to maintain activity that sensory cortex cannot. This picture does not obviously transfer to the mouse. Comparative anatomy finds no monotonic spine-density gradient in the mouse cortex, and available mouse-cortex connectomes~\cite{oh2014mesoscale,gamanut2018mouse,knox2018highresolution,harris2019hierarchical} indicate a connectivity structure that is denser and less hierarchical than in primates. The mechanism that allows the mouse cortex to maintain distributed working memory therefore cannot be inherited from the primate model.

In this paper we develop an anatomically-constrained large-scale model of the mouse cortex for working memory that does not rely on a gradient of recurrent excitation. The model is built on the mesoscopic Allen connectome~\cite{oh2014mesoscale,harris2019hierarchical,knox2018highresolution}, uses a mean-field reduction of a spiking attractor circuit~\cite{wang2002probabilistic,wongwang2006recurrent} to describe each of 43 cortical areas, and incorporates two measured cell-type features: a macroscopic gradient of parvalbumin (PV) interneuron density from the qBrain cell-type atlas~\cite{kim2017braintype}, and a hierarchy-dependent counterstream inhibitory bias (CIB) in which feedforward projections preferentially excite target excitatory cells and feedback projections preferentially recruit local inhibition~\cite{mejias2022mechanisms,wang2020macroscopic}. The model reproduces delay-period persistent activity in frontal and lateral association areas after transient stimulation of a primary sensory area, and can be extended with a 40-area thalamus to capture the thalamocortical support of mnemonic activity~\cite{guo2017maintenance,bolkan2017thalamic,schmitt2017thalamic}.

A central question is whether anatomical connectivity can predict the emergent large-scale activity pattern. Standard graph-theoretic summaries of the connectome that ignore cell-type targets cannot. We find that raw input strength and raw eigenvector centrality are essentially uncorrelated with delay-period firing rate across areas. To address this, we introduce cell type-specific graph measures that weight incoming projections by the cell-type composition of the target area and by CIB. These measures predict delay-period firing rate, predict whether each area sustains persistent activity, and when applied to reciprocal loops identify a small set of core areas whose inactivation collapses the distributed memory state. The same framework reveals that the same connectome can support many coexisting attractor states, with different tasks potentially engaging different cores. The contributions of this work are therefore: (i) an anatomically-grounded mouse-cortex working memory model that does not assume a gradient of recurrent excitation; (ii) the demonstration that inhibition (via PV gradient and CIB) jointly shapes distributed mnemonic activity and maintains an inhibitory stabilised baseline~\cite{tsodyks1997paradoxical,sanzeni2020inhibition}; (iii) cell type-specific connectivity measures that link the connectome to dynamics where classical measures fail; and (iv) the identification of a core subnetwork and a catalogue of coexisting attractors supported by the mesoscopic connectome.

\section{Related Work}

\paragraph{Multiregional models of cognition.}
Connectome-based large-scale cortical models have predominantly targeted the macaque~\cite{murray2014hierarchy,mejias2022mechanisms,froudistwalsh2021dopamine}. These models couple mean-field local circuits~\cite{wang2002probabilistic,wongwang2006recurrent} through tract-tracing-derived interareal weights and layer gradients of recurrent excitation or dopamine receptor density. They successfully reproduce a hierarchy of intrinsic timescales, graded persistent activity, and the prefrontal dominance of working memory codes, but they rely on macroscopic excitation gradients that are absent or weak in the mouse. Our model shares the same mean-field framework but replaces the excitation gradient with a measured PV inhibition gradient~\cite{kim2017braintype} and a cell-type-resolved CIB~\cite{mejias2022mechanisms,wang2020macroscopic}.

\paragraph{Mouse mesoscopic connectomics.}
The mouse cortical connectome has been mapped at mesoscopic resolution by several groups using viral tracing and retrograde tract tracing~\cite{oh2014mesoscale,gamanut2018mouse}, data-driven voxel models~\cite{knox2018highresolution}, and layer-resolved projections with hierarchy assignments~\cite{harris2019hierarchical}. Cell-type atlases now supply complementary information about the local composition of each area, including PV-interneuron density~\cite{kim2017braintype} and multimodal gradients of cytoarchitecture and transcription~\cite{fulcher2019overlapping}. These resources make it possible, for the first time, to build a mouse-cortex model that is simultaneously connectome-based and cell-type-aware; prior models treated interareal connectivity as unstructured by cell type.

\paragraph{Distributed mnemonic activity in mouse.}
Causal and recording studies in mice have repeatedly implicated a small set of frontal and thalamic regions in short-term memory and preparatory activity~\cite{guo2017maintenance,inagaki2019discrete,bolkan2017thalamic,schmitt2017thalamic,gilad2018behavioral}. In several of these studies, optogenetic silencing during the delay period collapses selectivity in otherwise-persistent areas, consistent with distributed maintenance rather than a single storage site. Our model-based inactivation analysis extends this body of work by asking, at the whole-cortex scale, which areas are causally required to maintain persistent activity and which merely inherit it. This yields a quantitative distinction between core, readout, input, and nonessential areas.

\paragraph{Inhibitory stabilisation and network stability.}
The idea that cortical networks operate in an inhibition-stabilised regime, where recurrent excitation alone would be unstable without fast inhibitory feedback, has strong theoretical and experimental support~\cite{tsodyks1997paradoxical,sanzeni2020inhibition}. Our large-scale model naturally inherits this property: removing either local inhibition or long-range projections onto inhibitory neurons destabilises the baseline state, while leaving them intact permits bistability of a subset of areas. The ISN framework therefore helps explain how distributed persistent activity can coexist with a stable low-firing baseline in the mouse cortex.

\paragraph{Graph-theoretic predictors of regional function.}
Structural connectivity has long been used to predict functional measures in the human brain, typically through input strength, degree, or eigenvector centrality. In our setting these raw measures fail to predict which areas maintain mnemonic activity. The gap between structure and function is closed by cell type-specific variants that incorporate PV density and hierarchy-dependent targeting, consistent with the broader argument~\cite{wang2020macroscopic,fulcher2019overlapping} that the functional relevance of interareal connections depends on which cell classes they target.
